\section{Planned high-dimensional and application benchmarks}
\label{app:benchmark_extensions}

The main experiments isolate basin-support alignment in controlled low-dimensional systems. The next benchmark additions should test the same claim in settings where autoencoding and regime discovery are more demanding, while preserving the paper's label-free training protocol. In both proposed benchmarks, labels are hidden from the encoder, decoder, transition model, sparsity rule, support-family construction, and checkpoint selection. Labels are used only after training to evaluate whether support families are useful regime variables.

\subsection{Spatialized multibasin reaction--diffusion fields}
\label{app:spatialized_multibasin_pde}

The low-risk high-dimensional benchmark is a spatialized version of the existing procedural multibasin generators. For a selected source system \(s\), let \(f_s:\mathbb{R}^2\to\mathbb{R}^2\) be the same two-dimensional vector field used in the main benchmark, with known attractor centers used only for evaluation. The spatial field is
\begin{equation}
  u(\xi,t)\in\mathbb{R}^2,\qquad \xi\in[0,1]^2,
\end{equation}
with dynamics
\begin{equation}
  \partial_t u(\xi,t)=f_s(u(\xi,t))+\nu\Delta u(\xi,t).
  \label{eq:spatialized_pde}
\end{equation}
On an \(N\times N\) periodic grid, the discretized dynamics are
\begin{equation}
  \dot u_{i,j}
  =
  f_s(u_{i,j})
  +
  \nu\left(
    u_{i+1,j}+u_{i-1,j}+u_{i,j+1}+u_{i,j-1}-4u_{i,j}
  \right),
  \label{eq:spatialized_grid}
\end{equation}
so the observed state dimension is \(2N^2\), e.g. \(8192\) for \(N=64\). The first pass should use \(N=64\), \(\nu\in\{0.005,0.01,0.02\}\), RK4 internal step size \(0.005\) or \(0.01\), periodic boundary conditions, and stored observations every 10--20 integration steps. The model observes only stored fields, not vector fields or integration substeps. This follows the storage and high-dimensional forecasting style of PDEBench and The Well, which provide reference conventions and baseline context for time-dependent PDE datasets \citep{takamoto_pdebench_2022,ohana_well_2024}.

Initial conditions should create spatially mixed local basin membership rather than independent ODE copies. For each trajectory, choose \(r\in\{2,3\}\) attractor centers \(c_1,\ldots,c_r\), sample low-frequency Fourier noise fields \(g_k(\xi)\), and convert them to soft masks
\begin{equation}
  w_k(\xi)=
  \frac{\exp(g_k(\xi)/\tau)}
       {\sum_{m=1}^r \exp(g_m(\xi)/\tau)}.
\end{equation}
The initial condition is
\begin{equation}
  u(\xi,0)=\sum_{k=1}^r w_k(\xi)c_k+\epsilon(\xi),
\end{equation}
where \(\epsilon(\xi)\) is small optionally smoothed Gaussian noise. Boundary-like cases can be oversampled by requiring each selected attractor to occupy a nontrivial initial area fraction. This produces fields with domain walls and coarsening-like transients.

Trajectory labels are computed after simulation and are never passed to the model. At a long label horizon \(T_{\rm label}\), each grid site is assigned to its nearest attractor center,
\begin{equation}
  b_{i,j}=\argmin_m \norm{u_{i,j}(T_{\rm label})-c_m}_2,
\end{equation}
and the trajectory-level final basin label is the modal site label \(B_{\rm global}\). The final majority fraction
\begin{equation}
  \rho_{\rm maj}=\max_m \frac{1}{N^2}\sum_{i,j}\mathbf{1}\{b_{i,j}=m\}
\end{equation}
records how much of the final field belongs to the modal basin. The primary support-alignment read should use all held-out trajectories, with \(B_{\rm global}\) hidden until evaluation, matching the main paper's all-held-out forecasting convention and avoiding an extra benchmark-specific filter. The majority fraction is a diagnostic stratification variable for asking whether support alignment is stronger on nearly single-basin fields than on mixed-domain fields; it is not a required filter for the main benchmark claim.

The initial dataset should cover five source systems, five seeds, and one grid size before expansion. Recommended source systems are \texttt{cal\_square\_4}, \texttt{cal\_high\_cross\_3}, \texttt{transition\_routes\_4}, \texttt{gated\_local\_linear}, and \texttt{var\_l\_shape\_5}; \texttt{cal\_octagon\_8}, \texttt{duffing\_triple\_well}, and \texttt{snic\_multi} are escalation cases. Store HDF5 arrays as \([{\rm trajectory},{\rm time},x,y,{\rm channel}]\), with metadata for source system, attractor centers, final basin label, majority fraction, per-pixel final basin map, stored \(\Delta t\), diffusion coefficient, and initial-condition seed.

Models should use convolutional encoders and decoders rather than flattening into the current MLP. The main comparison is a convolutional LISTA SKAE with dense \(K\), latent dimension \(d_z=512\), two or three LISTA refinement loops, and a convolutional decoder back to the two-channel field. Controls should include a dense Conv-KAE, a sparse Conv-MLP KAE without LISTA shrinkage, an FNO or U-Net autoregressive predictor, and an oracle basin-local Conv-KAE trained with evaluation labels only as an upper-bound diagnostic. Forecast metrics should include field MSE, gradient MSE, Fourier-shell MSE, and final-basin consistency. Support-family metrics should be computed out of sample: build \(F_{\rm abs}\) representatives on validation states only, freeze them, assign test states by nearest Jaccard representative, and report \(H(B_{\rm global}\mid F_{\rm abs})\), NMI, ARI, purity, \(|F_{\rm abs}|\), and \(H(F_{\rm abs}\mid B_{\rm global})\).

\subsection{Contact-rich robotic insertion in ManiSkill}
\label{app:maniskill_insertion}

The higher-risk application-style benchmark is robotic insertion in ManiSkill3 \citep{tao_maniskill3_2024}. The first task should be \texttt{PegInsertionSide-v1}, with \texttt{PlugCharger-v1} as a stretch task after the first signal is clear. These tasks contain randomized insertion geometry and demonstrations, making them suitable for testing whether sparse supports discover contact and outcome regimes rather than only synthetic basin labels.

Because the current KAE is autonomous, the robotics benchmark should use a controlled SKAE. Given observation \(o_t\) and action \(a_t\),
\begin{equation}
  z_t=\Enc(o_t),\qquad
  z_{t+1}^{\rm roll}=Kz_t+B\phi(a_t),\qquad
  \hat o_{t+1}=\Dec(z_{t+1}^{\rm roll}).
  \label{eq:controlled_skae}
\end{equation}
The support is still extracted only from \(z_t\), not from labels or hand-coded contact modes. The minimum-change alternative is to replay a fixed policy and model the closed-loop observation dynamics autonomously, but the controlled version is the cleaner extension because it separates action effects from the discovered support variable.

The first dataset should be state-only, not pixel-based. Start from downloaded demonstrations, recreate each environment from its metadata, replay the demonstrated action sequence, and generate perturbed rollouts using action noise and fixed action or end-effector biases. The rollout set should be deliberately balanced across successful insertions, rim jams, free-space misses, drops, and unstable partial insertions. Compact observations should include robot configuration and velocity, end-effector pose, gripper state, peg pose and velocity, box or receptacle pose, relative peg-hole pose, insertion depth, distance from peg tip to hole center, contact indicators when available, and the previous action. A first target scale is \(10{,}000\) rollouts, 100--150 steps each, split by reset seed and object geometry; after the state-only result works, RGB-D wrist and fixed-camera observations can be added.

Outcome and phase labels are evaluation-only. A trajectory-level outcome \(B\) can distinguish success, jammed-at-rim, dropped, missed-no-contact, and partial-insert-unstable outcomes using the simulator success flag plus contact persistence, peg-hole distance, insertion depth, gripper/peg relation, and final stability rules. A timestep-level contact phase \(C_t\) can label free-space approach, grasped transport, alignment, first contact, sliding insertion, jam, seated success, and drop or recovery. These labels are never used for training the model or constructing support families.

The robotics evaluation mirrors the basin-support protocol. Forecasting horizons should include 10, 25, 50, and 100 steps, with object-pose MSE, relative peg-hole pose MSE, insertion-depth error, contact-mode prediction F1, and outcome accuracy from the predicted rollout. Support-family evaluation should build \(F_{\rm abs}\) on validation states only and assign all held-out test states by frozen Jaccard representatives, reporting \(H(B\mid F_{\rm abs})\), NMI, ARI, \(H(C_t\mid F_{\rm abs})\), and contact-phase purity. The most important mechanistic figure is a support-family timeline over an insertion episode, aligned with camera frames, insertion depth, and contact signal. The intervention analogue is a success-to-jam or jam-to-success support swap at first contact under the same action sequence, measuring changes in predicted insertion depth, success probability, and contact rollout.

Baselines should include a dense controlled KAE, a sparse MLP controlled KAE without LISTA shrinkage, a GRU or Transformer world model, a mixture-of-linear-dynamics model with a learned soft gate, and a contact-oracle local linear model as an upper bound. The benchmark is worth promoting into the main evidence only if SKAE improves long-horizon relative-pose or insertion-depth prediction over the dense controlled KAE, support families align with outcome or contact labels better than dense and sparse-MLP controls, support swaps change predicted insertion outcomes under fixed actions, and the support timeline visibly changes at contact events.
